\documentclass[11pt]{article}
% The preamble every project's paper shares: % The preamble every project's paper shares: % The preamble every project's paper shares: \input{../../preamble} from
% projects/<project>/paper/. Boards come from chessboard and xskak, plots from
% pgfplots; every number a paper prints comes from its generated.tex.
\usepackage[LSBC4,T1]{fontenc}
\usepackage[utf8]{inputenc}
\usepackage{lmodern}
\usepackage[margin=1in]{geometry}
\usepackage{microtype}
\usepackage{amsmath}
\usepackage{booktabs}
\usepackage{tabularx}
\usepackage{array}
\usepackage{xcolor}
\usepackage[obeyspaces,spaces]{url}
\usepackage{xskak}
\usepackage{pgfplots}
\pgfplotsset{compat=1.18}
\usepackage{caption}
\usepackage{subcaption}
\usepackage{enumitem}
\usepackage[section]{placeins}
\usepackage[hidelinks]{hyperref}

\setchessboard{boardfontencoding=LSBC4, boardfontfamily=skaknew, blackfieldcolor=black!22,
  whitefieldcolor=white, setfontcolors, showmover=false, label=true, labelfontsize=6pt,
  boardfontsize=14pt}

% A move or line in UCI, and a position in FEN, both breakable at a slash.
\newcommand{\uci}[1]{\texttt{#1}}
\newcommand{\fen}[1]{{\footnotesize\path{#1}}}
\newcolumntype{L}{>{\raggedright\arraybackslash}X}

\setlength{\parskip}{0.4em}
\setlength{\parindent}{0pt}
 from
% projects/<project>/paper/. Boards come from chessboard and xskak, plots from
% pgfplots; every number a paper prints comes from its generated.tex.
\usepackage[LSBC4,T1]{fontenc}
\usepackage[utf8]{inputenc}
\usepackage{lmodern}
\usepackage[margin=1in]{geometry}
\usepackage{microtype}
\usepackage{amsmath}
\usepackage{booktabs}
\usepackage{tabularx}
\usepackage{array}
\usepackage{xcolor}
\usepackage[obeyspaces,spaces]{url}
\usepackage{xskak}
\usepackage{pgfplots}
\pgfplotsset{compat=1.18}
\usepackage{caption}
\usepackage{subcaption}
\usepackage{enumitem}
\usepackage[section]{placeins}
\usepackage[hidelinks]{hyperref}

\setchessboard{boardfontencoding=LSBC4, boardfontfamily=skaknew, blackfieldcolor=black!22,
  whitefieldcolor=white, setfontcolors, showmover=false, label=true, labelfontsize=6pt,
  boardfontsize=14pt}

% A move or line in UCI, and a position in FEN, both breakable at a slash.
\newcommand{\uci}[1]{\texttt{#1}}
\newcommand{\fen}[1]{{\footnotesize\path{#1}}}
\newcolumntype{L}{>{\raggedright\arraybackslash}X}

\setlength{\parskip}{0.4em}
\setlength{\parindent}{0pt}
 from
% projects/<project>/paper/. Boards come from chessboard and xskak, plots from
% pgfplots; every number a paper prints comes from its generated.tex.
\usepackage[LSBC4,T1]{fontenc}
\usepackage[utf8]{inputenc}
\usepackage{lmodern}
\usepackage[margin=1in]{geometry}
\usepackage{microtype}
\usepackage{amsmath}
\usepackage{booktabs}
\usepackage{tabularx}
\usepackage{array}
\usepackage{xcolor}
\usepackage[obeyspaces,spaces]{url}
\usepackage{xskak}
\usepackage{pgfplots}
\pgfplotsset{compat=1.18}
\usepackage{caption}
\usepackage{subcaption}
\usepackage{enumitem}
\usepackage[section]{placeins}
\usepackage[hidelinks]{hyperref}

\setchessboard{boardfontencoding=LSBC4, boardfontfamily=skaknew, blackfieldcolor=black!22,
  whitefieldcolor=white, setfontcolors, showmover=false, label=true, labelfontsize=6pt,
  boardfontsize=14pt}

% A move or line in UCI, and a position in FEN, both breakable at a slash.
\newcommand{\uci}[1]{\texttt{#1}}
\newcommand{\fen}[1]{{\footnotesize\path{#1}}}
\newcolumntype{L}{>{\raggedright\arraybackslash}X}

\setlength{\parskip}{0.4em}
\setlength{\parindent}{0pt}

% Written by projects/detector-themes/paper/generate.mjs; do not edit by hand.

\newcommand{\AppCommit}{d8b771238c}
\newcommand{\AppPuzzles}{4,630}
\newcommand{\HeldPuzzles}{20,000}
\newcommand{\TableBest}{%
\begin{tabular}{lrrrrrrrr}
\toprule
 & \multicolumn{4}{c}{the app's pool} & \multicolumn{4}{c}{held out} \\
\cmidrule(lr){2-5}\cmidrule(lr){6-9}
concept & tagged & said & precision & recall & tagged & said & precision & recall \\
\midrule
checkmate & 2074 & 1069 & 1.00 & 0.52 & 7833 & 3770 & 1.00 & 0.48 \\
fork & 563 & 389 & 0.92 & 0.63 & 2808 & 1804 & 0.91 & 0.58 \\
discovered-attack & 322 & 0 & \multicolumn{2}{c}{never said} & 1100 & 0 & \multicolumn{2}{c}{never said} \\
skewer & 196 & 0 & \multicolumn{2}{c}{never said} & 511 & 0 & \multicolumn{2}{c}{never said} \\
pin & 325 & 0 & \multicolumn{2}{c}{never said} & 1069 & 0 & \multicolumn{2}{c}{never said} \\
hanging-piece & 609 & 682 & 0.60 & 0.67 & 813 & 2025 & 0.28 & 0.70 \\
defensive-move & 315 & 864 & 0.03 & 0.09 & 637 & 5166 & 0.02 & 0.14 \\
back-rank-mate & 346 & 0 & \multicolumn{2}{c}{never said} & 845 & 12 & 0.00 & 0.00 \\
trapped-piece & 170 & 0 & \multicolumn{2}{c}{never said} & 209 & 0 & \multicolumn{2}{c}{never said} \\
deflection & 135 & 0 & \multicolumn{2}{c}{never said} & 789 & 0 & \multicolumn{2}{c}{never said} \\
attraction & 113 & 0 & \multicolumn{2}{c}{never said} & 640 & 0 & \multicolumn{2}{c}{never said} \\
f7-raid & 267 & 30 & 0.77 & 0.09 & 180 & 91 & 0.19 & 0.09 \\
double-check & 13 & 0 & \multicolumn{2}{c}{never said} & 80 & 0 & \multicolumn{2}{c}{never said} \\
zwischenzug & 34 & 0 & \multicolumn{2}{c}{never said} & 214 & 0 & \multicolumn{2}{c}{never said} \\
removing-the-defender & 5 & 0 & \multicolumn{2}{c}{never said} & 118 & 0 & \multicolumn{2}{c}{never said} \\
promotion & 361 & 0 & \multicolumn{2}{c}{never said} & 468 & 0 & \multicolumn{2}{c}{never said} \\
passed-pawn & 615 & 0 & \multicolumn{2}{c}{never said} & 1072 & 0 & \multicolumn{2}{c}{never said} \\
\bottomrule
\end{tabular}
}
\newcommand{\TableBlunder}{%
\begin{tabular}{lrrrrrrrr}
\toprule
 & \multicolumn{4}{c}{the app's pool} & \multicolumn{4}{c}{held out} \\
\cmidrule(lr){2-5}\cmidrule(lr){6-9}
concept & tagged & said & precision & recall & tagged & said & precision & recall \\
\midrule
fork & 563 & 262 & 0.94 & 0.44 & 2808 & 1374 & 0.92 & 0.45 \\
discovered-attack & 322 & 178 & 0.90 & 0.50 & 1100 & 683 & 0.79 & 0.49 \\
skewer & 196 & 0 & \multicolumn{2}{c}{never said} & 511 & 0 & \multicolumn{2}{c}{never said} \\
pin & 325 & 0 & \multicolumn{2}{c}{never said} & 1069 & 0 & \multicolumn{2}{c}{never said} \\
hanging-piece & 609 & 1619 & 0.38 & 1.00 & 813 & 6334 & 0.13 & 1.00 \\
defensive-move & 315 & 0 & \multicolumn{2}{c}{never said} & 637 & 0 & \multicolumn{2}{c}{never said} \\
back-rank-mate & 346 & 0 & \multicolumn{2}{c}{never said} & 845 & 0 & \multicolumn{2}{c}{never said} \\
trapped-piece & 170 & 0 & \multicolumn{2}{c}{never said} & 209 & 0 & \multicolumn{2}{c}{never said} \\
deflection & 135 & 61 & 0.46 & 0.21 & 789 & 548 & 0.55 & 0.38 \\
attraction & 113 & 146 & 0.18 & 0.24 & 640 & 883 & 0.23 & 0.32 \\
f7-raid & 267 & 0 & \multicolumn{2}{c}{never said} & 180 & 0 & \multicolumn{2}{c}{never said} \\
double-check & 13 & 0 & \multicolumn{2}{c}{never said} & 80 & 0 & \multicolumn{2}{c}{never said} \\
zwischenzug & 34 & 42 & 0.67 & 0.82 & 214 & 268 & 0.49 & 0.61 \\
removing-the-defender & 5 & 55 & 0.04 & 0.40 & 118 & 361 & 0.24 & 0.73 \\
promotion & 361 & 0 & \multicolumn{2}{c}{never said} & 468 & 0 & \multicolumn{2}{c}{never said} \\
passed-pawn & 615 & 0 & \multicolumn{2}{c}{never said} & 1072 & 0 & \multicolumn{2}{c}{never said} \\
\bottomrule
\end{tabular}
}
\newcommand{\TableMove}{%
\begin{tabular}{lrrrrrrrr}
\toprule
 & \multicolumn{4}{c}{the app's pool} & \multicolumn{4}{c}{held out} \\
\cmidrule(lr){2-5}\cmidrule(lr){6-9}
concept & tagged & said & precision & recall & tagged & said & precision & recall \\
\midrule
checkmate & 2074 & 0 & \multicolumn{2}{c}{never said} & 7833 & 0 & \multicolumn{2}{c}{never said} \\
fork & 563 & 387 & 0.92 & 0.63 & 2808 & 1802 & 0.91 & 0.58 \\
discovered-attack & 322 & 246 & 0.79 & 0.61 & 1100 & 979 & 0.66 & 0.59 \\
skewer & 196 & 153 & 0.91 & 0.71 & 511 & 446 & 0.65 & 0.57 \\
pin & 325 & 111 & 0.73 & 0.25 & 1069 & 472 & 0.54 & 0.24 \\
hanging-piece & 609 & 0 & \multicolumn{2}{c}{never said} & 813 & 0 & \multicolumn{2}{c}{never said} \\
defensive-move & 315 & 0 & \multicolumn{2}{c}{never said} & 637 & 0 & \multicolumn{2}{c}{never said} \\
back-rank-mate & 346 & 263 & 0.21 & 0.16 & 845 & 774 & 0.18 & 0.17 \\
trapped-piece & 170 & 200 & 0.63 & 0.74 & 209 & 494 & 0.32 & 0.76 \\
deflection & 135 & 0 & \multicolumn{2}{c}{never said} & 789 & 0 & \multicolumn{2}{c}{never said} \\
attraction & 113 & 0 & \multicolumn{2}{c}{never said} & 640 & 0 & \multicolumn{2}{c}{never said} \\
f7-raid & 267 & 14 & 1.00 & 0.05 & 180 & 8 & 1.00 & 0.04 \\
double-check & 13 & 5 & 1.00 & 0.38 & 80 & 43 & 1.00 & 0.54 \\
zwischenzug & 34 & 0 & \multicolumn{2}{c}{never said} & 214 & 0 & \multicolumn{2}{c}{never said} \\
removing-the-defender & 5 & 0 & \multicolumn{2}{c}{never said} & 118 & 0 & \multicolumn{2}{c}{never said} \\
promotion & 361 & 16 & 1.00 & 0.04 & 468 & 34 & 0.97 & 0.07 \\
passed-pawn & 615 & 454 & 0.73 & 0.54 & 1072 & 806 & 0.64 & 0.48 \\
\bottomrule
\end{tabular}
}
\newcommand{\DefSaidApp}{864}
\newcommand{\DefSaidHeld}{5,166}
\newcommand{\DefPrecApp}{0.03}
\newcommand{\DefPrecHeld}{0.02}
\newcommand{\DefShareHeld}{26}
\newcommand{\PhaseShareApp}{34}
\newcommand{\PhaseShareHeld}{36}
\newcommand{\HangSaidApp}{1619}
\newcommand{\HangRealApp}{609}
\newcommand{\HangSaidHeld}{6,334}
\newcommand{\HangRealHeld}{813}
\newcommand{\HangPrecApp}{0.38}
\newcommand{\HangPrecHeld}{0.13}
\newcommand{\BestHangPrecApp}{0.60}
\newcommand{\BestHangPrecHeld}{0.28}
\newcommand{\BrmSaidHeld}{12}
\newcommand{\BrmTpHeld}{0}
\newcommand{\BrmSaidApp}{0}
\newcommand{\MovePinRecHeld}{0.24}
\newcommand{\MoveBrmPrecHeld}{0.18}
\newcommand{\NamedMates}{\texttt{pillsburysMate}, \texttt{operaMate}, \texttt{epauletteMate}, \texttt{smotheredMate}, \texttt{cornerMate}, \texttt{swallowstailMate}, \texttt{morphysMate}, \texttt{hookMate}, \texttt{arabianMate}, \texttt{killBoxMate}, \texttt{triangleMate}, \texttt{anastasiaMate}, \texttt{blindSwineMate}, \texttt{bodenMate}, \texttt{dovetailMate}, \texttt{doubleBishopMate}, \texttt{vukovicMate}, \texttt{balestraMate}}
\newcommand{\NamedMateTags}{1,051}
\newcommand{\PinBoard}{%
\chessboard[setfen={8/5n2/p1K1k3/P7/5P2/4b3/8/3R4 w - - 3 43}, boardfontsize=13pt, pgfstyle=straightmove, arrow=to, linewidth=0.07em, color=red!75!black, markmoves={d4-e3}, pgfstyle=straightmove, arrow=to, linewidth=0.07em, color=green!50!black, markmoves={d1-e1}]
}
\newcommand{\PinId}{0MHcM}
\newcommand{\PinRating}{1462}
\newcommand{\PinFen}{\fen{8/5n2/p1K1k3/P7/5P2/4b3/8/3R4 w - - 3 43}}
\newcommand{\PinFirst}{42\ldots{}\symbishop{}e3}
\newcommand{\PinLine}{43.~\symrook{}e1 \symknight{}e5+ 44.~fxe5}
\newcommand{\PinUci}{\uci{d1e1 f7e5 f4e5}}
\newcommand{\PinThemes}{\texttt{crushing}, \texttt{endgame}, \texttt{pin}, \texttt{short}}
\newcommand{\PinBest}{\texttt{defensiveMove}}
\newcommand{\PinBlunder}{nothing}
\newcommand{\PinMove}{\texttt{trapped-piece}, \texttt{pin}, \texttt{open-file}}
\newcommand{\SkewerBoard}{%
\chessboard[setfen={6rk/2R2K2/1p2p3/4P3/1P6/4B3/6p1/8 b - - 1 41}, boardfontsize=13pt, pgfstyle=straightmove, arrow=to, linewidth=0.07em, color=red!75!black, markmoves={h6-e3}, pgfstyle=straightmove, arrow=to, linewidth=0.07em, color=green!50!black, markmoves={g8-g7}]
}
\newcommand{\SkewerId}{1dno2}
\newcommand{\SkewerRating}{1178}
\newcommand{\SkewerFen}{\fen{6rk/2R2K2/1p2p3/4P3/1P6/4B3/6p1/8 b - - 1 41}}
\newcommand{\SkewerFirst}{41.~\symbishop{}e3}
\newcommand{\SkewerLine}{41\ldots{}\symrook{}g7+ 42.~\symking{}f6 \symrook{}xc7}
\newcommand{\SkewerUci}{\uci{g8g7 f7f6 g7c7}}
\newcommand{\SkewerThemes}{\texttt{crushing}, \texttt{endgame}, \texttt{short}, \texttt{skewer}}
\newcommand{\SkewerBest}{\texttt{defensiveMove}}
\newcommand{\SkewerBlunder}{nothing}
\newcommand{\SkewerMove}{\texttt{skewer}}
\newcommand{\ForkBoard}{%
\chessboard[setfen={4rn2/p2bkpR1/2p1p3/1p2P2B/5n2/1PN5/P2K1P2/8 w - - 0 35}, boardfontsize=13pt, pgfstyle=straightmove, arrow=to, linewidth=0.07em, color=red!75!black, markmoves={d5-f4}, pgfstyle=straightmove, arrow=to, linewidth=0.07em, color=green!50!black, markmoves={g7-f7}]
}
\newcommand{\ForkId}{1Aofp}
\newcommand{\ForkRating}{1168}
\newcommand{\ForkFen}{\fen{4rn2/p2bkpR1/2p1p3/1p2P2B/5n2/1PN5/P2K1P2/8 w - - 0 35}}
\newcommand{\ForkFirst}{34\ldots{}\symknight{}xf4}
\newcommand{\ForkLine}{35.~\symrook{}xf7+ \symking{}d8 36.~\symrook{}xf4}
\newcommand{\ForkUci}{\uci{g7f7 e7d8 f7f4}}
\newcommand{\ForkThemes}{\texttt{crushing}, \texttt{fork}, \texttt{middlegame}, \texttt{short}}
\newcommand{\ForkBest}{\texttt{fork}}
\newcommand{\ForkBlunder}{\texttt{hangingPiece}}
\newcommand{\ForkMove}{\texttt{fork}, \texttt{king-attack}}
\newcommand{\CandAppAgreeBase}{40.7}
\newcommand{\CandAppAgree}{49.0}
\newcommand{\CandAppAgreeDiff}{8.3}
\newcommand{\CandAppAgreeCi}{[7.5, 9.1]}
\newcommand{\CandAppDisDiff}{$-$2.0}
\newcommand{\CandAppDisCi}{[$-$2.5, $-$1.5]}
\newcommand{\CandHeldAgreeBase}{30.3}
\newcommand{\CandHeldAgree}{35.9}
\newcommand{\CandHeldAgreeDiff}{5.6}
\newcommand{\CandHeldAgreeCi}{[5.2, 5.9]}
\newcommand{\CandHeldDisDiff}{$-$0.2}
\newcommand{\CandHeldDisCi}{[$-$0.5, 0.1]}
\newcommand{\CandDraws}{2,000}
\newcommand{\DeckMoved}{44}
\newcommand{\DeckMovedFit}{20}
\newcommand{\DeckPins}{24}
\newcommand{\DeckPinsFit}{8}

% Written by projects/detector-themes/paper/generate.mjs --round2; do not edit by hand.

\newcommand{\RTwoAppCommit}{f4bd63f19b343f686fb97ea4cf70b211c0bddc0a}
\newcommand{\RTwoFrozenCommit}{b055ee132c0cde90c2409f52ae9a94c8d9afa0a6}
\newcommand{\RTwoPuzzles}{4,630}
\newcommand{\RTwoSampled}{20,000}
\newcommand{\RTwoEligible}{4,682,966}
\newcommand{\RTwoRatingLow}{400}
\newcommand{\RTwoRatingHigh}{1,899}
\newcommand{\RTwoPoolHash}{005f31596fae53c51d76916faca5de4b02850d2beaea693e21e99ff574b40f85}
\newcommand{\RTwoSampleHash}{98c6777e6dbc7b81e0a7763d8005f3c12eb5ccc01da41a208103fa036614037a}
\newcommand{\RTwoDatabaseHash}{76335bfa7d7c4a7f93c1366d81549e53951ebb79dd43d34904cab8f22d962f8d}
\newcommand{\RTwoSeed}{20261014}
\newcommand{\RTwoBootstrapSeed}{20261014}
\newcommand{\RTwoDraws}{2,000}
\newcommand{\RTwoConfidence}{95}
\newcommand{\RTwoDisagreementLimit}{1}
\newcommand{\RTwoCounts}{%
\begin{tabular}{llrr}
\toprule
reader & set & agrees & disagrees \\
\midrule
pinned & app pool & 2,269 & 1,058 \\
pinned & held out & 7,179 & 6,759 \\
deflection & app pool & 2,308 & 1,044 \\
deflection & held out & 7,526 & 6,712 \\
trapped piece & app pool & 2,376 & 1,044 \\
trapped piece & held out & 7,324 & 6,769 \\
\bottomrule
\end{tabular}
}
\newcommand{\RTwoTable}{%
\begin{tabular}{llrrl}
\toprule
candidate & set & agreement change & disagreement change & rule \\
\midrule
deflection & app pool & 0.84 [0.583, 1.102] & $-$0.30 [$-$0.562, $-$0.043] & pass \\
deflection & held out & 1.74 [1.565, 1.920] & $-$0.24 [$-$0.420, $-$0.055] & pass \\
trapped piece & app pool & 2.31 [1.901, 2.765] & $-$0.30 [$-$0.454, $-$0.151] & pass \\
trapped piece & held out & 0.73 [0.615, 0.850] & 0.05 [$-$0.020, 0.115] & pass \\
\bottomrule
\end{tabular}
}
\newcommand{\RTwoVerdicts}{\texttt{best-deflection}: keep; \texttt{best-trapped}: keep}


\title{Do the app's motif labels name the tactic a position holds?\\[0.3em]
\large Three of the coaching app's board readers, scored on public puzzles}
\author{Chess Coach research lab}
\date{}

\begin{document}
\maketitle

\noindent The main text records the earlier study at commit \texttt{\AppCommit}.
Its results and examples are preserved. The appendix reports the later app-pool
development phase at a different pin; it contains no new held-out result.

\begin{abstract}
Three of the app's board readers put a motif word on a position: the theme of the best move the
player missed (it files the misses deck, the puzzle buckets and the rule ledger, and the learner
credits a solve to the skill it names), the motifs a bad move allowed (counted by the weakness
profile, the metrics and the pattern census), and what a move does (the explanations' source).
We scored each against the themes Lichess's own tagger put on \AppPuzzles{} puzzles the app ships
and on \HeldPuzzles{} held-out puzzles of the same ratings. All three name forks well. The misses
deck's labeller has no word for pins, skewers, discovered attacks, deflections or trapped
pieces, and when it says ``defensive move'' (\DefShareHeld\% of held-out puzzles) the tagger
agrees \DefPrecHeld{} of the time. The review scan tags a hanging piece on \HangSaidHeld{}
held-out puzzles, against \HangRealHeld{} that the tagger says have one. The move describer is
the most precise of the three and the one to build the others on: naming its skewers, discovered
attacks and pins before the misses deck's defensive fallback raised agreement on both sets
(\CandAppAgreeDiff{} and \CandHeldAgreeDiff{} points) and is kept.
\end{abstract}

\section{The question}

A motif word is a claim about the board. When the misses deck files a position under
``defensive move'', the learner counts a solve there toward a defensive-move skill; when the
review scan tags a hanging piece, the weakness profile counts one more hanging piece. A wrong
word sends practice and advice to the wrong place, so before the routines lean on these words
we asked how often each is right when it is said (precision) and how often it is said when it
is right (recall).

A Lichess puzzle is a mistake, its first move, followed by the line that punishes it, and
Lichess tags each with themes. That is the shape the app's readers already read: the position
the solver faces, the best move, the mistake that allowed it and the refutation. Lichess's tags
come from another heuristic tagger, not from the truth, so this measures agreement; the three
puzzles in Section~\ref{sec:boards} read disagreements on the board.

\section{Method}

Each reader runs on each puzzle the way the app runs it on a game (commit
\texttt{\AppCommit}): the theme of the best move on the solver's position and the first
solution move (\texttt{themeOfBest}); the motifs of the mistake from the position before it,
the mistake, the solution line and the score it reaches (\texttt{blunderMotifs}); and the
findings that are a gain for the solver when the first solution move is described with the rest
of the line as its reply (\texttt{describeMove}). Both sides' words are mapped to concepts by the
app's own vocabulary, so a Lichess theme and an app word name one concept the way the app
resolves them. A concept that only one side has a word for is not scored. The app's pool is the
training set for any change tuned here; the held-out set, a seeded uniform sample of the same
rating range that leaves out every puzzle the app ships, is the test.

\section{Results}

\begin{table}[tbp]
\centering
\small
\TableBest
\caption{The misses deck's labeller (\texttt{themeOfBest}). Tagged: puzzles Lichess tags with
the concept; said: puzzles the reader names it on. Checkmate is said only for a mate in one.}
\end{table}

\begin{table}[tbp]
\centering
\small
\TableBlunder
\caption{The review scan's motifs (\texttt{blunderMotifs}). Checkmate is left out: the score
passed in decides it.}
\end{table}

\begin{table}[tbp]
\centering
\small
\TableMove
\caption{The move describer's gain findings (\texttt{describeMove}).}
\end{table}

\paragraph{The misses deck's labeller.} It has six motif words. When none fires it files the
position under its phase (\PhaseShareApp\% of the app's pool, \PhaseShareHeld\% held out) or, when
one of the solver's pieces was loose and the move protected it, as a defensive move. Puzzles
are tactics, so the defensive move mostly catches the tactics the labeller has no word for:
said on \DefSaidApp{} of the app's pool and \DefSaidHeld{} held-out puzzles, it agrees with the
tagger \DefPrecApp{} and \DefPrecHeld{} of the time. Its hanging piece is right \BestHangPrecApp{}
and \BestHangPrecHeld{} of the time. It said back-rank mate on \BrmSaidApp{} of the app's pool
and \BrmSaidHeld{} held-out puzzles, \BrmTpHeld{} of them tagged so. Its fork is good.

\paragraph{The review scan.} It tags a hanging piece whenever the refutation wins material, so
it tags one on \HangSaidApp{} puzzles of the app's pool (\HangRealApp{} tagged) and \HangSaidHeld{}
held out (\HangRealHeld{} tagged): precision \HangPrecApp{} and \HangPrecHeld. It never tags a
pin, a skewer, a back-rank mate or a trapped piece. Its fork and discovered attack are precise
but catch about half.

\paragraph{The move describer.} The most precise of the three on the tactics both
vocabularies name: fork, skewer and discovered attack all above 0.6 precision and 0.5 recall on
both sets. It finds a quarter of pins (\MovePinRecHeld{} held out), and its back-rank mate is
mostly not one the tagger sees (\MoveBrmPrecHeld).

\paragraph{The vocabulary.} The app's vocabulary does not resolve the named mates Lichess
tags (\NamedMates; \NamedMateTags{} tags held out), so they score as nothing on either side.

\section{Three held-out puzzles}\label{sec:boards}

Each board is the position the solver faces; the red arrow is the mistake just played, the
green one the first solution move.

\begin{figure}[h]
\centering
\begin{minipage}[t]{0.31\textwidth}\centering
\PinBoard\\[0.3em]
\footnotesize Puzzle \texttt{\PinId}, rated \PinRating
\end{minipage}\hfill
\begin{minipage}[t]{0.31\textwidth}\centering
\SkewerBoard\\[0.3em]
\footnotesize Puzzle \texttt{\SkewerId}, rated \SkewerRating
\end{minipage}\hfill
\begin{minipage}[t]{0.31\textwidth}\centering
\ForkBoard\\[0.3em]
\footnotesize Puzzle \texttt{\ForkId}, rated \ForkRating
\end{minipage}
\end{figure}

\paragraph{\PinId.} After \PinFirst, the solution \PinLine{} (\PinUci) from \PinFen. Lichess:
\PinThemes. The misses deck's labeller: \PinBest. The describer: \PinMove.

\paragraph{\SkewerId.} After \SkewerFirst, the solution \SkewerLine{} (\SkewerUci) from
\SkewerFen. Lichess: \SkewerThemes. The misses deck's labeller: \SkewerBest. The describer:
\SkewerMove.

\paragraph{\ForkId.} After \ForkFirst, the solution \ForkLine{} (\ForkUci) from \ForkFen.
Lichess: \ForkThemes. The review scan: \ForkBlunder. The misses deck's labeller: \ForkBest.

\section{A first candidate}

The rule for a change, fixed before the first one ran: per puzzle, a reader \emph{agrees} when one
of its words names a concept the puzzle is tagged with and \emph{disagrees} when one names a
concept it is not tagged with; a change is kept when, on both sets, agreement rises with the 95\%
interval of the paired difference (\CandDraws{} bootstrap draws of puzzles) above zero and
disagreement does not rise.

The first candidate changes only the misses deck's labeller: before its defensive fallback, it
names a skewer, a discovered attack or a pin when the move describer finds one for the mover.
Agreement rose from \CandAppAgreeBase\% to \CandAppAgree\% on the app's pool (\CandAppAgreeDiff{}
points, \CandAppAgreeCi) and from \CandHeldAgreeBase\% to \CandHeldAgree\% held out
(\CandHeldAgreeDiff{} points, \CandHeldAgreeCi); disagreement moved \CandAppDisDiff{} points
\CandAppDisCi{} and \CandHeldDisDiff{} points \CandHeldDisCi. It is kept.

On the player's misses deck the same change re-files \DeckMoved{} positions, and the best line
wins material or mates in \DeckMovedFit{} of them; the describer's pin, on \DeckPins{} positions,
wins material in \DeckPinsFit, because it also fires on pins that win nothing at once. A port
should carry the skewer and the discovered attack, and the pin only when the line wins with it.

\section{What it means for the app}

\begin{enumerate}[nosep]
\item Port the kept candidate: the describer's skewer and discovered attack before the
defensive fallback, and its pin only when the line wins with it; then make the defensive move
require that the move answers a threat the board code can name, the next candidate.
\item The review scan's hanging piece should require a piece that was attacked and not
defended before the refutation, not any material the refutation wins; its counts feed the
weakness profile and the metrics.
\item The vocabulary should resolve the named mates, so a mate pattern is a concept on both
sides.
\item Any change is tuned on the app's pool and kept only if it holds on the held-out set.
\end{enumerate}

\section{Limits}

Lichess's tags are another tagger's, and a puzzle has as many tags as that tagger found, so a
correct app word with no matching tag counts as wrong. Puzzles are tactics by construction; the
misses deck is built from the player's own mistakes, where a defensive move is more common, which
the drill-decks project measures on the deck itself.

\appendix
\section{App-pool development at the later pin}

This phase measures commit \texttt{\RTwoAppCommit} on the \RTwoPuzzles{} unique
puzzles shipped by that app. Each candidate starts from that pin and changes one
labeller. The app's existing readers supply the board claims. Candidates are
developed on this pool, finalized, and committed before the held-out test.

The new sample contains \RTwoSampled{} puzzles drawn uniformly from
\RTwoEligible{} eligible database rows, rated \RTwoRatingLow{} through
\RTwoRatingHigh{}, with the app pool excluded and seed \RTwoSeed. Only its sample
metadata is read here. The app-pool data hash is
\texttt{\RTwoPoolHash}; the sample hash is \texttt{\RTwoSampleHash}, and its
database hash is \texttt{\RTwoDatabaseHash}.

\begin{table}[htbp]
\centering
\small
\RTwoTable
\caption{Development counts on the app pool. Agrees and disagrees count puzzles;
changed counts a different best-move word. The last columns count puzzles where
the candidate's new word matches a tag or does not. These are development counts,
not a keep verdict.}
\end{table}

Agreement increases on the app pool for every finalized candidate in the table.
That is only the development screen. The keep rule remains unchanged: on both
the app pool and held-out set, the paired agreement difference must have its
\RTwoConfidence\% interval wholly above zero, and the disagreement difference's
upper endpoint must be at most \RTwoDisagreementLimit{} percentage point.
The test uses \RTwoDraws{} bootstrap draws of puzzles, with replacement, and
seed \RTwoBootstrapSeed{} restarted for each set. There is no multiple-comparison
correction in this fixed rule. No held-out labeller or candidate comparison was
run in this phase, so none of these candidates is yet kept.

The app-pool work summed to \RTwoAppSeconds{} seconds of recorded per-puzzle
wall time across the baseline and candidates. Scaling that work to the sample
at \RTwoProcesses{} processes estimates \RTwoHeldSeconds{} seconds for scoring.
This excludes process launch, candidate copies and comparison. The local sample
pass took \RTwoSampleSeconds{} seconds; downloading its database is separate.
The estimate is a planning figure, not a timed held-out run.

\end{document}
